\documentclass[11pt]{article}
\usepackage[utf8]{inputenc}
\usepackage{amsmath,amssymb}
\usepackage[numbers,sort&compress]{natbib}
\usepackage{hyperref}

\title{Certified Lagrangian lower bounds from local-solver multipliers
in a rigorous interval branch and bound}
\author{Draft --- outline}
\date{November 2026}

\begin{document}
\maketitle

\begin{abstract}
\textbf{[Draft abstract.]} Interval branch-and-bound solvers such as IbexOpt compute
rigorous enclosures of the global optimum of nonconvex continuous problems, but their
lower bounds come from linear relaxations of the problem over each box and stay loose
long after the incumbent is optimal. Meanwhile, the local NLP solver that these codes
call for upper bounding returns, at no extra cost, a KKT point together with Lagrange
multipliers that are discarded. We show that this by-product is enough to compute a
rigorous lower bound on the original nonconvex problem at every node of the search: we enclose the Hessian of the Lagrangian with the solver's multipliers by
interval arithmetic and either certify that it is positive semidefinite on the box, or
shift it diagonally in the spirit of $\alpha$BB. In both cases the tangent plane of the
Lagrangian at the local solution, minimized in closed form over the box and corrected for
the relaxation of equality constraints, yields a valid lower bound. In IbexOpt with Ipopt
as local solver, on 298 instances of the COCONUT benchmark, the bound closes 8 more instances and
roughly halves the solving time on the instances that take more than a second, at a
cost below 1\% of the running time, and makes the rigorous solver close more instances
than SCIP~10 on this benchmark.
\end{abstract}

\section{Introduction}

We consider the nonconvex continuous optimization problem
\begin{equation}\label{eq:P}
  \min_{x \in X} \; f(x) \quad \text{s.t.} \quad g(x) \le 0, \quad h(x) = 0,
\end{equation}
where $X \subset \mathbb{R}^n$ is a box and $f$, $g$ and $h$ are twice
differentiable functions given by explicit expressions. \emph{Rigorous} global
optimization solvers return an interval that provably contains the global minimum value,
even in the presence of rounding errors. Interval branch-and-bound codes such as
GlobSol \citep{Kearfott1996}, IBBA \citep{Ninin2015} and IbexOpt
\citep{Trombettoni2011,Araya2014} reach this guarantee by evaluating every quantity
that affects pruning with outward-rounded interval arithmetic
\citep{HansenWalster2004,Neumaier2004}. The equalities are relaxed to
$|h(x)| \le \varepsilon_h$, and the solver certifies feasible points inside the
relaxed region.

In such a solver, each node of the search tree is a box $[x] \subseteq X$ on which the
solver (i)~contracts the box with constraint propagation and with linear relaxations,
(ii)~computes a lower bound of $f$ over the feasible part of the box, typically as the
optimal value of a linear relaxation made safe with the technique of
\citet{NeumaierShcherbina2004} and \citet{Jansson2004lp}, and (iii)~searches for a
feasible point that may improve the incumbent value (the \emph{loup}). In IbexOpt,
step~(iii) combines cheap heuristics (probing and the inner regions of
\citet{Araya2014}) with calls to a local NLP solver, the interior-point solver Ipopt
\citep{WachterBiegler2006}, whose approximate local optimum is then certified as
feasible by a small box built around it. Calling a local solver for upper bounding is
standard practice in global optimization; nothing in what follows is specific to Ipopt
beyond the fact that it returns Lagrange multipliers.

Two observations motivate this work. First, on a large benchmark the incumbent is usually
optimal, or very close to it, long before the search ends. Most of the remaining effort
goes into closing the gap from below, and the linear relaxations of the problem over
small boxes are the bottleneck. Second, each successful call to the local solver returns not only a
point $\hat x$ but a full primal--dual pair $(\hat x, \hat\mu, \hat\lambda)$, with
multipliers $\hat\mu \ge 0$ for the inequalities and $\hat\lambda$ for the equalities.
Interval solvers discard the multipliers. Yet weak Lagrangian duality holds for
\emph{any} multipliers, not only optimal ones: for every feasible $x$ of the relaxed
problem,
\begin{equation}\label{eq:weak}
  f(x) \;\ge\; L(x) - \varepsilon_h \textstyle\sum_j |\hat\lambda_j|,
  \qquad L(x) := f(x) + \hat\mu^\top g(x) + \hat\lambda^\top h(x).
\end{equation}
A lower bound on $f$ over the feasible part of $[x]$ therefore follows from any lower
bound of the Lagrangian $L$ over the \emph{whole} box, with no constraints left. When
$L$ is convex on $[x]$, its tangent plane at $\hat x$ is such a bound, and it is tight
when $\hat x$ is a KKT point of the problem restricted to $[x]$. This is the theorem of
\citet{Jansson2004convex} and \citet{MessineTrombettoni2019} for convex programs.

Our point is that the convexity hypothesis is needed only for $L$ with the
\emph{fixed} multipliers $(\hat\mu,\hat\lambda)$ and only \emph{on the current box}. It
can therefore be checked numerically, node by node, in the original nonconvex problem.
Let $[H]$ be an interval enclosure of $\nabla^2 L$ over $[x]$, and let $d \ge 0$ be a
vector such that $H + \operatorname{diag}(d)$ is positive semidefinite for every
$H \in [H]$. Following the scaled Gerschgorin rule of \citet{Adjiman1998a}, we take
$d_i = \max\{0, \sum_{j \ne i} \lvert [H]_{ij} \rvert - h^L_{ii}\}$, where $h^L_{ii}$ is the lower end of $[H]_{ii}$.
Then, for every $x \in [x]$,
\begin{equation}\label{eq:bound}
  L(x) \ge L(\hat x) + \nabla L(\hat x)^\top (x - \hat x)
          - \frac{1}{2} \sum_i d_i (x_i - \hat x_i)^2 .
\end{equation}
The right-hand side is separable, so it is bounded below over $[x]$ in closed form, in
$O(n)$ operations once the gradient is known. Combined with \eqref{eq:weak} and
evaluated in interval arithmetic, it yields a rigorous lower bound on~\eqref{eq:P}
over $[x]$, which raises the node's lower bound or prunes the node. When the Hessian
is certified positive semidefinite on the box ($d = 0$), the bound reduces to that of
\citet{Jansson2004convex}. It then converges to the local optimum value as the box
shrinks around a nondegenerate KKT point.

\paragraph{Contributions.}
\begin{enumerate}
\item \textbf{A certified Lagrangian lower bound for the original nonconvex problem}
  (main contribution). It reuses the multipliers of the local solver that already
  performs upper bounding, so it needs no relaxation, no subgradient iterations and no
  auxiliary optimization problem. Its only extra cost is one interval evaluation of
  the Hessian of the Lagrangian per call. It handles nonlinear equalities, through the
  $\varepsilon_h$-correction in~\eqref{eq:weak}, and boxes on which the Lagrangian is
  not convex, through the diagonal shift in~\eqref{eq:bound}; on boxes where convexity
  is certified it is the bound of \citet{Jansson2004convex}.
\item \textbf{Making the bound cheap}: a backoff that applies it only where it pays
  off, and positive-semidefiniteness tests of increasing strength (interval Gerschgorin,
  comparison matrix, midpoint--radius test with an interval Cholesky factorization).
  The bound costs less than 1\% of the running time and the time per node grows by at
  most 5\%.
\item \textbf{An experimental study} on 298 instances of the COCONUT benchmark \citep{Shcherbina2003}: an
  ablation (IbexOpt without local solver, with the local solver for upper bounding
  only, and with the Lagrangian bound) and a comparison with SCIP~10. With the
  Lagrangian bound, the rigorous solver closes 236 instances, against 228 without it
  and 220 for SCIP, and its solving time drops to 0.55 of the time without it on the
  instances that take more than a second.
  \textbf{[TODO: SCIP with $\pm 10^{12}$ bounds for both solvers.]}
\end{enumerate}

The paper is organized as follows. Section~\ref{sec:related} reviews related work.
Section~\ref{sec:prelim} recalls the interval branch and bound of IbexOpt and how it
uses a local solver. Section~\ref{sec:bound} presents the certified Lagrangian bound,
its rigorous evaluation, the convexity certificates and the diagonal shift, and its
integration in the search. Section~\ref{sec:exp} reports the experiments and
Section~\ref{sec:concl} concludes.

\section{Related work}\label{sec:related}

\subsection{Rigorous global optimization and upper bounding with local solvers}

Spatial branch-and-bound solvers for nonconvex problems include BARON
\citep{RyooSahinidis1996,TawarmalaniSahinidis2005}, Couenne \citep{Belotti2009},
ANTIGONE \citep{MisenerFloudas2014} and SCIP \citep{VigerskeGleixner2018}. They build
convex (mostly polyhedral) relaxations from a factorable reformulation and call a local
NLP solver to find feasible points. Their bounds are computed in floating point, with
tolerances, and are not guaranteed against rounding errors.

Interval solvers trade speed for rigor
\citep{Kearfott1996,HansenWalster2004,Neumaier2004}. In IbexOpt
\citep{ChabertJaulin2009,Trombettoni2011,Araya2014}, the lower bound comes from
interval evaluation and from linear relaxations, either XTaylor linearizations or
affine arithmetic \citep{Ninin2015,Araya2024hybrid}. These relaxations are made safe a
posteriori by the dual-based correction of
\citet{NeumaierShcherbina2004} \citep[see also][]{Jansson2004lp,Guilbeau}.

Feasible points are certified either by inner regions, i.e.\ polytopes or boxes proved
entirely feasible \citep{Araya2014}, or by an existence proof in a small box around an
approximate local optimum \citep{Kearfott2014}. Ipopt \citep{WachterBiegler2006} is
the local solver used by IbexOpt for the latter; we keep this upper-bounding
scheme unchanged and reuse its by-products.

How often a global solver calls its local solver is a design choice of every code:
BARON runs local searches in preprocessing and during the tree under its own rules,
Couenne solves on average a fixed number of local NLPs per tree level, and SCIP's
subNLP heuristic decides at each node from the nodes processed since its last run and
its past success \citep{VigerskeGleixner2018}; multistart is also common
\citep{Ugray2007}. In all of them, the local solver serves the upper bound only. Inside
rigorous solvers, local search has also been used for node selection
\citep{Neveu2016,ReyesAraya2024}, and Charibde \citep{Vanaret2020} shares the
incumbent between an evolutionary algorithm and an interval branch and bound.

\subsection{Rigorous bounds from approximate dual information}

Weak duality turns any dual-feasible vector into a valid bound. Making this rigorous
with approximate multipliers is a well-established idea for linear programs.

\citet{NeumaierShcherbina2004} and, independently, \citet{Jansson2004lp} showed how to
compute a safe lower bound on the optimal value of a linear program from an approximate
dual solution, by bounding the dual infeasibility over the finite simple bounds with
directed rounding. \citet{Guilbeau} review these bounds and correct typographical errors
in the original formulas. This correction is what makes linear relaxations safe inside
interval solvers, IbexOpt included.

\citet{Jansson2004convex} extended the approach to convex programs
$\min f(x)$ s.t.\ $G(x) \le 0$, $Hx = h$ with simple bounds, where $f$ and every
component of $G$ are convex and the equalities are linear. Each function is replaced by
its tangent at an approximate solution $\tilde x$. The dual of the resulting linear
program, evaluated at the approximate multipliers, gives the bound, which is exact at a
KKT point. The method is a post-processing of any nonlinear solver's output, and it
handles uncertain data and infinite bounds.
\citet{MessineTrombettoni2019} derived the same bound from a second-order Taylor
expansion of the Lagrangian and the KKT conditions, in the form
$f(x^*) \ge L(\tilde x) + (z - \tilde x)^\top \nabla L(\tilde x)$, with $z_i$ the end
of $[x_i]$ selected by the sign of $\partial_i L(\tilde x)$. They proposed it to make
convex relaxations solved by a local solver reliable inside interval codes such as Ibex
and IBBA.

QIBEX \citep{Elloumi2025} put this bound to work. It solves mixed-integer quadratically
constrained programs in IbexOpt with a convex quadratic relaxation obtained once from a
semidefinite program. It certifies both the convexity of the relaxation and, via the
same theorem, its lower bound.

All of these works apply the bound to a problem that is convex \emph{by construction}:
a linear program, a convex program, or a convex relaxation. Our bound applies it to the
original nonconvex problem. In the proofs of \citet{Jansson2004convex} and
\citet{MessineTrombettoni2019}, convexity of $f$ and $G$ and linearity of the equalities
serve only one purpose: making the Lagrangian with the \emph{given} multipliers convex
on the domain. We certify exactly this weaker condition, box by box, from an interval
Hessian. When it fails, we restore validity with a diagonal shift. The primal--dual pair
comes from the local solver already used for upper bounding, so no relaxation is built
or solved.

\subsection{Convex underestimators and the $\alpha$BB shift}

The $\alpha$BB method \citep{Androulakis1995,Adjiman1998a,Adjiman1998b} underestimates
each general nonconvex term $f$ on a box by
$f(x) + \sum_i \alpha_i (x_i^L - x_i)(x_i^U - x_i)$. This underestimator is convex if
and only if $\nabla^2 f + 2\operatorname{diag}(\alpha)$ is positive semidefinite on the
box \citep[Thm.~2.1]{Adjiman1998a}.

\citet{Adjiman1998a} give several rigorous ways to obtain $\alpha$ from an interval
Hessian:
\begin{itemize}
\item $O(n^2)$ rules: Gerschgorin and scaled Gerschgorin;
\item $O(n^3)$ rules: E-matrix, Mori--Kokame, the lower bounding Hessian and the Hertz
  method;
\item Kharitonov's theorem applied to the interval characteristic polynomial;
\item an SDP-based rule.
\end{itemize}
In the scaled Gerschgorin rule (their Thm.~3.13), for any $d > 0$, taking
$\alpha_i = \max\{0, -\frac{1}{2}(a^L_{ii} - \sum_{j\ne i} \lvert a \rvert_{ij} d_j/d_i)\}$
makes $A + 2\operatorname{diag}(\alpha)$ positive semidefinite for every $A \in [A]$.
With $d = \mathbf 1$ this is our shift; they also propose
$d = x^U - x^L$.

In their experiments \citep{Adjiman1998b}:
\begin{itemize}
\item computing $\alpha$ takes about 2.5\% of the CPU time for most rules;
\item the scaled Gerschgorin rule with $d = x^U - x^L$, the E-matrix rule and the Hertz
  rule are usually the best, but the differences are modest;
\item branching on the variables that most degrade the underestimator speeds up
  convergence.
\end{itemize}

Our use of the shift differs in two ways. First, we shift the \emph{aggregated
Lagrangian} with fixed multipliers rather than each nonconvex term, so cancellations
between objective and constraints, which can make $L$ convex on a box where $f$ and
$g$ are not, come for free. Second, $\alpha$BB solves a convex nonlinear program per
node with the underestimators. We use the shifted Lagrangian only through its tangent
at the local solution~\eqref{eq:bound}, a separable expression bounded in closed form.
The bound is weaker than the optimal value of the $\alpha$BB relaxation, but it costs
one interval Hessian evaluation instead of a nonlinear solve.

\subsection{Lagrangian underestimates in branch and bound}

\citet{VanVoorhis2002} is the closest predecessor: it uses the Lagrangian, not each
function, as the object to underestimate. He solves quadratically constrained programs
by a branch and bound in which the lower bound minimizes a convex underestimator of the
Lagrangian over the box. The underestimator is built term by term for quadratic
functions; the multipliers are updated by subgradient steps in each node; and interval
Newton on the Fritz John conditions discards boxes. That work does not address rounding
errors.

He observes that obtaining good multipliers takes several iterations, and suggests
using Lagrangian underestimates ``at later stages of the branch-and-bound process, when
good multiplier estimates have been obtained via RLT or some other relaxation
strategy''. Our approach takes the multipliers, at no cost, from the local solver that
already runs at those stages. It replaces the per-node convex program by the
closed-form bound~\eqref{eq:bound}, covers general twice-differentiable functions, and
is rigorous.

In BARON, the multipliers of the relaxation are used for optimality-based domain
reduction \citep{RyooSahinidis1996}. Extending our bound to reduce variable ranges
with the same tangent is a natural next step.

\subsection{Certifying convexity}

Deciding whether every matrix in a symmetric interval matrix is positive semidefinite
is NP-hard, so rigorous solvers rely on sufficient conditions. Examples include
Gerschgorin-type and comparison-matrix tests, midpoint--radius conditions such as
$\lambda_{\min}(H_c) \ge \rho(H_\Delta)$, and their parametric refinements
\citep{Hladik2018}.

Symbolic approaches propagate convexity through the expression graph
\citep{Fourer2010}. \citet{Kearfott2013} identify the (possibly intermediate)
variables responsible for nonconvexity, aiming at a non-adaptive subdivision of only
those variables. In interval branch and bound, the reverse use of the interval Hessian
is classical: the nonconvexity test discards boxes on which the objective cannot have
an interior minimizer \citep{HansenWalster2004}.

These analyses address the convexity of the \emph{problem functions}. Ours concerns the
Lagrangian with fixed multipliers on one box, and it is numerical and local. A symbolic
analysis could complement it by predicting on which instances the certification is
likely to succeed.


\section{Preliminaries}\label{sec:prelim}

\subsection{Interval branch and bound}

\textbf{[Draft.]} Intervals, boxes and outward rounding; an inclusion function $[f]$
of $f$ satisfies $f([x]) \subseteq [f]([x])$ for every box. IbexOpt keeps the
objective as an extra variable $y$ with the constraint $y = f(x)$, so that every node
is a box $[x] \times [y]$ and its lower bound is $\underline{y}$. Equalities are
relaxed to $|h_j(x)| \le \varepsilon_h$ ($\varepsilon_h = 10^{-8}$ by default) and
a range constraint $h_j(x) \in [a_j, b_j]$ (thick constant) is kept as such. The search
stops when $(\text{loup} - \text{uplo})$ is below a relative ($10^{-3}$) or absolute
($10^{-7}$) precision.

At each node, IbexOpt
(i)~contracts the box with constraint propagation (HC4, ACID) and with a linear
relaxation (XTaylor and affine arithmetic, \texttt{--relax both}) solved by an LP
solver whose optimum is made safe as in \citet{NeumaierShcherbina2004};
(ii)~raises $\underline{y}$ accordingly;
(iii)~runs the upper-bounding heuristics;
(iv)~stores the node in a best-first heap (\texttt{lsmear} branching).
\textbf{[TODO: one figure with the node loop and where the new bound plugs in.]}

\subsection{Upper bounding with a local solver}\label{sec:ub}

The loup finder combines probing, inner-region heuristics on the linear relaxation
\citep{Araya2014} and a local NLP solver (Ipopt) called every $k$ loup-finder calls
(and when another heuristic improves the incumbent, from its point). Ipopt solves the
problem restricted to the current box from its midpoint and returns
$(\hat x, \hat\mu, \hat\lambda)$. Because $\hat x$ satisfies the constraints only up to
Ipopt's tolerances, a short nested search in a small box around $\hat x$ looks for a
point that interval arithmetic proves feasible; only a certified point updates the
loup. We keep this scheme unchanged; what we add is to store the last usable
primal--dual pair $(\hat x, \hat\mu, \hat\lambda)$ that Ipopt returns.

\section{A certified Lagrangian lower bound}\label{sec:bound}

\subsection{Weak duality with relaxed equalities}

Write the constraints as $g(x) \le 0$ and $h(x) = 0$ (each $\ge$ constraint negated).
For any $\mu \ge 0$ and $\lambda$, and any $x$ that the search treats as feasible
($g(x) \le 0$, $|h(x)| \le \varepsilon_h$),
\[
  f(x) \;\ge\; L_{\mu,\lambda}(x) - \varepsilon_h \|\lambda\|_1,
  \qquad L_{\mu,\lambda} = f + \mu^\top g + \lambda^\top h .
\]
Ipopt's multipliers are projected onto the right signs (a multiplier of an
inequality with the wrong sign is set to $0$), which keeps the inequality valid without
any assumption on their quality.

\subsection{Tangent bound on a box}

\textbf{Proposition 1.} Let $[x]$ be a box, $\hat x \in [x]$, and $d \in \mathbb{R}^n_+$
such that $\nabla^2 L(\xi) + \operatorname{diag}(d) \succeq 0$ for all $\xi \in [x]$.
Then for all $x \in [x]$ feasible for the relaxed problem,
\[
  f(x) \ge L(\hat x) + \sum_i \min_{x_i \in [x_i]} \partial_i L(\hat x)(x_i - \hat x_i)
  - \frac12 \sum_i d_i \max_{x_i \in [x_i]} (x_i - \hat x_i)^2
  - \varepsilon_h \|\lambda\|_1 .
\]
\emph{Proof sketch.} $\phi(x) = L(x) + \frac12 \sum_i d_i (x_i - \hat x_i)^2$ is convex on
$[x]$ and has gradient $\nabla L(\hat x)$ at $\hat x$; the tangent inequality of $\phi$
and the bounds on each term give the result. \textbf{[TODO: full proof; remark that
with $d = 0$ it is the bound of \citet{Jansson2004convex} and
\citet{MessineTrombettoni2019}, and that the separable penalty could use the exact
minimum of each one-dimensional concave term over its interval.]}

\textbf{Exactness.} If $d = 0$, $\hat x$ is a KKT point of the problem restricted to
$[x]$ with multipliers $(\hat\mu, \hat\lambda)$ and computations are exact, the bound
equals $f(\hat x)$ up to $\varepsilon_h\|\lambda\|_1$ \citep[Thm.~2]{Jansson2004convex}.
With $d \neq 0$ the penalty is $O(\|d\|\,\mathrm{wid}([x])^2)$ and vanishes as the box
shrinks. \textbf{[TODO: state as a proposition on convergence along a nested sequence of
boxes.]}

\subsection{Rigorous evaluation}

All quantities are evaluated with interval arithmetic at the point interval
$\hat x$ (projected onto the node box): $L(\hat x)$, $\nabla L(\hat x)$, the products
$\partial_i L(\hat x)\,([x_i] - \hat x_i)$, whose lower bound is the minimum of the
linear term, and the penalty with the upper bound of $([x_i] - \hat x_i)^2$. The bound
is the lower end of the resulting interval. Multipliers are floating-point numbers
used as exact constants. Unbounded boxes, unbounded Hessian entries or non-finite
values give no bound ($-\infty$).

\subsection{Convexity certificates and the diagonal shift}

$[H]$ is an interval enclosure of $\nabla^2 L$ over $[x]$, obtained by automatic
differentiation of the gradient enclosures of $f$ and of each $g_k$, $h_k$ weighted by
the fixed multipliers. Two modes:
\begin{itemize}
\item \emph{certify}: prove $[H] \succeq 0$ (every symmetric $H \in [H]$), with tests of
increasing cost: interval Gerschgorin; the comparison matrix; the midpoint--radius
condition checked by an interval Cholesky factorization of $H_c - \rho I$.
\textbf{[TODO: precise statements, cost $O(n^2)$ / $O(n^3)$, reference to
\citet{Hladik2018}.]} If none succeeds, no bound.
\item \emph{shift} (default): $d_i = \max\{0, \sum_{j\ne i} |[H]_{ij}| - h^L_{ii}\}$,
the Gerschgorin rule of \citet[Thm.~3.13 with $d = \mathbf 1$]{Adjiman1998a}; the
scaled variant (weights $w_i = \max |[x_i] - \hat x_i|$) is also available.
\end{itemize}

\subsection{Integration in the search}

The bound is computed in the contraction step of every node, after the other
contractors, with the last stored primal--dual pair (computed at another node, so
nothing is solved at the current one). If it exceeds $\underline{y}$, it raises it; if
it exceeds the loup, the node is pruned. \textbf{Backoff:} when the bound does not
improve the node's lower bound, the next attempt is postponed by twice the previous
gap (up to 64 nodes); a useful attempt resets the gap. This bounds the cost on
problems where the Lagrangian bound never helps.
\textbf{[TODO: pseudocode.]}

\section{Experiments}\label{sec:exp}

\subsection{Setup}

\textbf{[Draft.]} 298 instances of the COCONUT benchmark \citep{Shcherbina2003}, in the Minibex
translation distributed with IbEX (\texttt{benchs/optim/all}): the 223 that IbexOpt solves within the time limit and
75 that it does not. One thread, 310~s, same machine, three runs in parallel; relative
precision $10^{-3}$, absolute $10^{-7}$, $\varepsilon_h = 10^{-8}$. Deterministic budgets
(backoff in nodes), seed 1. Measures: instances solved; time ratios with a 1~s floor
(timeout $=$ 310~s) on the instances some configuration solves and some takes more
than 1~s, with bootstrap confidence intervals; nodes; time per node.
\textbf{[TODO: hardware, Ipopt version, LP solver.]}

\subsection{Ablation}

Configurations: no local solver; local solver for upper bounding only (IbexOpt's
default); + Lagrangian bound with the shift (default); + Lagrangian bound with
certification only. Table: solved /298, time ratio [IC95], nodes, time per node,
share of time in the local solver and in the bound. Current results: 223, 228, 236
solved; time 1.62, 1.00, 0.55. \textbf{[TODO: certify-only arm; table.]}

\subsection{Where the bound helps}

How often the Lagrangian is certified convex on the node box, how often the shifted
bound improves the node bound or prunes it, as a function of depth; the instances it
closes (abel, abelbis, abelbis-1, abelter-1, catena, ex8\_4\_1, fccu, polak3) and the
ones it accelerates (robot, harker, sambal, dnieper\ldots); where it never helps.
\textbf{[TODO: statistics from \texttt{--cvx-stats}; one figure: fraction of useful
nodes vs depth.]}

\subsection{Variants of the bound}

Certify vs shift; plain vs scaled Gerschgorin; with and without backoff; PSD tests.
\textbf{[TODO: table on the instances where the bound is used.]}

\subsection{Comparison with SCIP}

SCIP~10 with the same stopping rule, both solvers on the same domains (missing bounds
set to $\pm 10^{12}$); pairwise performance profiles; solved instances per model class;
discussion of rigor (SCIP's feasibility tolerance; wrong optima reported in the
literature). \textbf{[TODO: results of the $\pm 10^{12}$ runs.]}

\section{Conclusion}\label{sec:concl}

\textbf{[TODO.]} Summary; limits (needs a usable local solution; the shift degrades on
wide boxes; multipliers from another node); future work: optimality-based range
reduction with the same tangent \citep{RyooSahinidis1996}, better shifts (scaled,
E-matrix), branching on the variables with $d_i > 0$ \citep{Adjiman1998b}.

\bibliographystyle{plainnat}
\bibliography{refs}
\end{document}
